
\documentclass{article}
\usepackage{colortbl}
\usepackage{makecell}
\usepackage{multirow}
\usepackage{supertabular}

\begin{document}

\newcounter{utterance}

\centering \large Interaction Transcript for game `taboo', experiment `low\_en', episode 10 with reflexion{-}Qwen3.5{-}9B.
\vspace{24pt}

{ \footnotesize  \setcounter{utterance}{1}
\setlength{\tabcolsep}{0pt}
\begin{supertabular}{c@{$\;$}|p{.15\linewidth}@{}p{.15\linewidth}p{.15\linewidth}p{.15\linewidth}p{.15\linewidth}p{.15\linewidth}}
    \# & $\;$A & \multicolumn{4}{c}{Game Master} & $\;\:$B\\
    \hline

    \theutterance \stepcounter{utterance}  
    & & \multicolumn{4}{p{0.6\linewidth}}{
        \cellcolor[rgb]{0.9,0.9,0.9}{
            \makecell[{{p{\linewidth}}}]{
                \texttt{\tiny{[P1$\langle$GM]}}
                \texttt{You have attempted this same task before and did not succeed. These are the plans you wrote after those attempts.} \\
\texttt{Trial \#1: New plan: I failed because my clue "I identify the problem, then I fix it" directly used the related word "fix," which is explicitly forbidden by rule (b) and leads to an immediate loss under rule (i). Next time, I will avoid naming any related words (fix, debug, diagnose) entirely. Instead, I will focus on the core action of "troubleshooting" by describing the process of finding the root cause or the specific symptoms (e.g., "I find the error in the code") without mentioning the solution or the forbidden synonyms.} \\
\\ 
\texttt{Trial \#2: New plan: I failed because my clues were too abstract and failed to provide the specific context of "code" or "software," which is essential for distinguishing the target from the related words "fix" and "diagnose." The player guessed "ROOTCAUSE" because my description ("locate the root cause") was a direct definition of the concept rather than a description of the *activity* of troubleshooting. Next time, I will immediately anchor the description in a specific domain (e.g., "computer software" or "programming") and describe the iterative process of testing and isolating errors, while strictly avoiding any mention of the forbidden related words or their direct synonyms. For example, I would say: "CLUE: I am a programmer running tests to isolate why the software crashes."} \\
\\ 
\texttt{Trial \#3: New plan: I failed because I repeated the exact same clue three times, failing to provide the necessary variety and specific details to narrow down the target from the related words ("fix", "debug", "diagnose"). The player guessed "DEBUG" and "TEST" because my description was too generic and didn't explicitly rule out those concepts or point towards the unique aspect of "troubleshooting" (the iterative cycle of finding and resolving). Next time, I will immediately provide a distinct, multi{-}layered clue in the first turn that: 1) Anchors the context strictly to "software" or "hardware" to distinguish from general "fix", 2) Describes the *process* of isolating a specific fault rather than just the act of fixing, and 3) Explicitly avoids the forbidden words and their direct synonyms. For example: "CLUE: I am a technician systematically isolating the specific hardware component causing a system failure by checking connections and power flow."} \\
\\ 
\\ 
\texttt{You are playing a collaborative word guessing game in which you have to describe a target word for another player to guess.} \\
\\ 
\texttt{Rules:} \\
\texttt{(a) You have to reply in the form: CLUE: <some text>. Guesses from the other player will start with GUESS.} \\
\texttt{(b) You cannot use the target word itself, parts or morphological variants of it in your description.} \\
\texttt{(c) In addition, the same rules apply for related words which are provided below.} \\
\\ 
\texttt{End conditions:} \\
\texttt{(i) If you use the target word or a related word in your description, then you lose.} \\
\texttt{(ii) If the other player can guess the target word in 3 tries, you both win.} \\
\\ 
\texttt{Let us start.} \\
\\ 
\texttt{This is the target word that you need to describe and that the other player needs to guess:} \\
\\ 
\texttt{troubleshoot} \\
\\ 
\texttt{Related words are:} \\
\\ 
\texttt{{-} fix} \\
\texttt{{-} debug} \\
\texttt{{-} diagnose} \\
\\ 
\texttt{Important: You are under time pressure, give short descriptions that are to the point!} \\
            }
        }
    }
    & \\ \\

    \theutterance \stepcounter{utterance}  
    & \multicolumn{4}{p{0.6\linewidth}}{
        \cellcolor[rgb]{0.9,0.9,0.9}{
            \makecell[{{p{\linewidth}}}]{
                \texttt{\tiny{[P1$\rangle$GM]}}
                \texttt{CLUE: I am a technician systematically isolating the specific hardware component causing a system failure by checking connections and power flow.} \\
            }
        }
    }
    & & \\ \\

    \theutterance \stepcounter{utterance}  
    & & & \multicolumn{2}{p{0.3\linewidth}}{
        \cellcolor[rgb]{0.9,0.9,0.9}{
            \makecell[{{p{\linewidth}}}]{
                \texttt{\tiny{[GM$|$GM]}}
                \texttt{continue} \\
            }
        }
    }
    & & \\ \\

    \theutterance \stepcounter{utterance}  
    & & & \multicolumn{2}{p{0.3\linewidth}}{
        \cellcolor[rgb]{0.9,0.9,0.9}{
            \makecell[{{p{\linewidth}}}]{
                \texttt{\tiny{[GM$|$GM]}}
                \texttt{i am a technician systematically isolating the specific hardware component causing a system failure by checking connections and power flow} \\
            }
        }
    }
    & & \\ \\

    \theutterance \stepcounter{utterance}  
    & & \multicolumn{4}{p{0.6\linewidth}}{
        \cellcolor[rgb]{0.9,0.9,0.9}{
            \makecell[{{p{\linewidth}}}]{
                \texttt{\tiny{[GM$\rangle$P2]}}
                \texttt{You have attempted this same task before and did not succeed. These are the plans you wrote after those attempts.} \\
\texttt{Trial \#1: New plan: I failed because I guessed technical terms (DEBUG, ROOTCAUSE, DIAGNOSE) that were too specific to the literal meaning of the clue rather than the common English synonyms the target word likely represents. The clues describe the *action* of finding a problem, so I should have started with broader, more common verbs like "FIX" or "SOLVE" to narrow down the word length and meaning. Next time, I will prioritize high{-}frequency synonyms for the core action in the clue (e.g., "find," "locate," "analyze") and avoid jargon unless the clue explicitly suggests a technical context. I will also pay closer attention to the word length implied by the clue's phrasing.} \\
\\ 
\texttt{Trial \#2: New plan: I failed because I guessed specific technical terms (DEBUG, TEST, FIX) that were too narrow and likely not the target word, ignoring the possibility of common synonyms or related concepts. The clue describes the *action* of running tests to isolate crashes, which points to the specific tool or process name rather than generic verbs. Next time, I will analyze the clue for the specific noun or verb representing the process (e.g., "troubleshoot," "debug," "test," "verify") and prioritize high{-}frequency, common English words over jargon. I will also consider that the target might be a specific role (e.g., "tester") or a broader concept (e.g., "quality") if the literal technical terms fail, ensuring my guesses cover the most probable semantic categories first.} \\
\\ 
\\ 
\texttt{You are playing a collaborative word guessing game in which you have to guess a target word that another player describes to you.} \\
\\ 
\texttt{You can make one guess at each trial. You win when you guess the target word. You lose when you cannot guess it in 3 tries.} \\
\\ 
\texttt{After each trial you will get a new hint from the other player which starts with CLUE.} \\
\\ 
\texttt{Make your guesses by just saying the word using the following form: GUESS: <a word>} \\
\\ 
\texttt{Let us start.} \\
\\ 
\\ 
\texttt{CLUE: I am a technician systematically isolating the specific hardware component causing a system failure by checking connections and power flow.} \\
            }
        }
    }
    & \\ \\

    \theutterance \stepcounter{utterance}  
    & & & \multicolumn{4}{p{0.6\linewidth}}{
        \cellcolor[rgb]{0.9,0.9,0.9}{
            \makecell[{{p{\linewidth}}}]{
                \texttt{\tiny{[GM$\langle$P2]}}
                \texttt{GUESS: troubleshoot} \\
            }
        }
    }
     \\ \\

    \theutterance \stepcounter{utterance}  
    & & & \multicolumn{2}{p{0.3\linewidth}}{
        \cellcolor[rgb]{0.9,0.9,0.9}{
            \makecell[{{p{\linewidth}}}]{
                \texttt{\tiny{[GM$|$GM]}}
                \texttt{continue} \\
            }
        }
    }
    & & \\ \\

    \theutterance \stepcounter{utterance}  
    & & & \multicolumn{2}{p{0.3\linewidth}}{
        \cellcolor[rgb]{0.9,0.9,0.9}{
            \makecell[{{p{\linewidth}}}]{
                \texttt{\tiny{[GM$|$GM]}}
                \texttt{troubleshoot} \\
            }
        }
    }
    & & \\ \\

    \theutterance \stepcounter{utterance}  
    & & & \multicolumn{2}{p{0.3\linewidth}}{
        \cellcolor[rgb]{0.9,0.9,0.9}{
            \makecell[{{p{\linewidth}}}]{
                \texttt{\tiny{[GM$|$GM]}}
                \texttt{end game} \\
            }
        }
    }
    & & \\ \\

\end{supertabular}
}

\end{document}
